
\documentclass{article}
\usepackage{colortbl}
\usepackage{makecell}
\usepackage{multirow}
\usepackage{supertabular}

\begin{document}

\newcounter{utterance}

\centering \large Interaction Transcript for game `codenames', experiment `opponent\_easy', episode 1 with claude{-}opus{-}4{-}8{-}azure.
\vspace{24pt}

{ \footnotesize  \setcounter{utterance}{1}
\setlength{\tabcolsep}{0pt}
\begin{supertabular}{c@{$\;$}|p{.15\linewidth}@{}p{.15\linewidth}p{.15\linewidth}p{.15\linewidth}p{.15\linewidth}p{.15\linewidth}}
    \# & $\;$A & \multicolumn{4}{c}{Game Master} & $\;\:$B\\
    \hline

    \theutterance \stepcounter{utterance}  
    & & & \multicolumn{2}{p{0.3\linewidth}}{
        \cellcolor[rgb]{0.9,0.9,0.9}{
            \makecell[{{p{\linewidth}}}]{
                \texttt{\tiny{[GM$|$GM]}}
                {'hidden': {'team': ['bermuda', 'smrt', 'kanada', 'netopir', 'london', 'miška', 'stadion', 'spojina', 'princesa'], 'innocent': ['nebotičnik', 'trikotnik', 'strežnik', 'helikopter', 'antarktika', 'zob', 'indija'], 'opponent': ['diamant', 'zaloga', 'krt', 'pajek', 'sesalnik', 'marmor', 'skrinja', 'feniks'], 'assassin': ['vlak']}, 'revealed': {'team': {'team': [], 'innocent': [], 'opponent': [], 'assassin': []}, 'opponent': {'team': [], 'innocent': [], 'opponent': [], 'assassin': []}}}
            }
        }
    }
    & & \\ \\

    \theutterance \stepcounter{utterance}  
    & & \multicolumn{4}{p{0.6\linewidth}}{
        \cellcolor[rgb]{0.9,0.9,0.9}{
            \makecell[{{p{\linewidth}}}]{
                \texttt{\tiny{[P1$\langle$GM]}}
                \texttt{Navedite eno enobesedno iztočnico, ki se nanaša na nekatere besede vaše ekipe s spodnjega seznama. Izberete lahko, na koliko in na katere besede se nanaša vaša iztočnica.} \\
\texttt{Cilj je, da vaš partner na podlagi vaših iztočnic ugane vse besede vaše ekipe, preden druga ekipa odkrije vse svoje besede. Če vaš partner ugane besedo "nasprotnika" ali "motečo" besedo, se njegovo ugibanje konča in na vrsto pride nasprotna ekipa. Če vaš partner ugane besedo "ubij", se igra konča in vaša ekipa takoj izgubi.} \\
\texttt{Več besed, na katere se nanaša vaša iztočnica, pomeni hitrejšo zmago v igri, vendar bo vaš partner morda težje uganil asociacije, zato pametno izberite iztočnico. Beseda iztočnica mora biti pomensko povezana s ciljnimi besedami, ne sme biti ena od besed na seznamih ali vsebovati njihovih delov.} \\
\\ 
\texttt{Vedno navedite enobesedno iztočnico in z vejico ločen seznam sorodnih ciljnih besed v naslednji obliki, odgovori naj bodo čim krajši, nikoli ne vključite drugega besedila, kot je zahtevano v tej obliki:} \\
\\ 
\texttt{NAMIG: <BESEDA>} \\
\texttt{CILJI: <CILJI>} \\
\\ 
\texttt{Vaše ekipne besede so: bermuda, smrt, kanada, netopir, london, miška, stadion, spojina, princesa.} \\
\texttt{Besede nasprotnikove ekipe so: diamant, zaloga, krt, pajek, sesalnik, marmor, skrinja, feniks.} \\
\texttt{Besede, ki odvračajo pozornost, so: nebotičnik, trikotnik, strežnik, helikopter, antarktika, zob, indija.} \\
\texttt{Besede za ubijanje so: vlak.} \\
            }
        }
    }
    & \\ \\

    \theutterance \stepcounter{utterance}  
    & \multicolumn{4}{p{0.6\linewidth}}{
        \cellcolor[rgb]{0.9,0.9,0.9}{
            \makecell[{{p{\linewidth}}}]{
                \texttt{\tiny{[P1$\rangle$GM]}}
                \texttt{NAMIG: geografija} \\
\texttt{CILJI: bermuda, kanada, london} \\
            }
        }
    }
    & & \\ \\

    \theutterance \stepcounter{utterance}  
    & & & \multicolumn{2}{p{0.3\linewidth}}{
        \cellcolor[rgb]{0.9,0.9,0.9}{
            \makecell[{{p{\linewidth}}}]{
                \texttt{\tiny{[GM$|$GM]}}
                \texttt{geografija} \\
            }
        }
    }
    & & \\ \\

    \theutterance \stepcounter{utterance}  
    & & & \multicolumn{2}{p{0.3\linewidth}}{
        \cellcolor[rgb]{0.9,0.9,0.9}{
            \makecell[{{p{\linewidth}}}]{
                \texttt{\tiny{[GM$|$GM]}}
                ['bermuda', 'kanada', 'london']
            }
        }
    }
    & & \\ \\

    \theutterance \stepcounter{utterance}  
    & & & \multicolumn{2}{p{0.3\linewidth}}{
        \cellcolor[rgb]{0.9,0.9,0.9}{
            \makecell[{{p{\linewidth}}}]{
                \texttt{\tiny{[GM$|$GM]}}
                {'word': 'bermuda', 'assignment': 'team'}
            }
        }
    }
    & & \\ \\

    \theutterance \stepcounter{utterance}  
    & & & \multicolumn{2}{p{0.3\linewidth}}{
        \cellcolor[rgb]{0.9,0.9,0.9}{
            \makecell[{{p{\linewidth}}}]{
                \texttt{\tiny{[GM$|$GM]}}
                {'word': 'kanada', 'assignment': 'team'}
            }
        }
    }
    & & \\ \\

    \theutterance \stepcounter{utterance}  
    & & & \multicolumn{2}{p{0.3\linewidth}}{
        \cellcolor[rgb]{0.9,0.9,0.9}{
            \makecell[{{p{\linewidth}}}]{
                \texttt{\tiny{[GM$|$GM]}}
                {'word': 'london', 'assignment': 'team'}
            }
        }
    }
    & & \\ \\

    \theutterance \stepcounter{utterance}  
    & & \multicolumn{4}{p{0.6\linewidth}}{
        \cellcolor[rgb]{0.9,0.9,0.9}{
            \makecell[{{p{\linewidth}}}]{
                \texttt{\tiny{[GM$\rangle$P2]}}
                \texttt{Navedite z vejico ločen seznam do 3 besed s spodnjega seznama, ki so najbolj povezane ali najbolj tesno povezane z besedo 'geografija'. Seznam ugibanj vedno začnite z UGIBANJE:  in v odgovor ne vključite nobenega drugega besedila.} \\
\\ 
\texttt{strežnik, helikopter, netopir, stadion, vlak, sesalnik, bermuda, diamant, nebotičnik, marmor, trikotnik, spojina, krt, zob, smrt, kanada, london, indija, miška, skrinja, pajek, antarktika, zaloga, feniks, princesa} \\
            }
        }
    }
    & \\ \\

    \theutterance \stepcounter{utterance}  
    & & & \multicolumn{4}{p{0.6\linewidth}}{
        \cellcolor[rgb]{0.9,0.9,0.9}{
            \makecell[{{p{\linewidth}}}]{
                \texttt{\tiny{[GM$\langle$P2]}}
                \texttt{UGIBANJE: kanada, indija, antarktika} \\
            }
        }
    }
     \\ \\

    \theutterance \stepcounter{utterance}  
    & & & \multicolumn{2}{p{0.3\linewidth}}{
        \cellcolor[rgb]{0.9,0.9,0.9}{
            \makecell[{{p{\linewidth}}}]{
                \texttt{\tiny{[GM$|$GM]}}
                ['kanada', 'indija', 'antarktika']
            }
        }
    }
    & & \\ \\

    \theutterance \stepcounter{utterance}  
    & & & \multicolumn{2}{p{0.3\linewidth}}{
        \cellcolor[rgb]{0.9,0.9,0.9}{
            \makecell[{{p{\linewidth}}}]{
                \texttt{\tiny{[GM$|$GM]}}
                {'word': 'kanada', 'assignment': 'team'}
            }
        }
    }
    & & \\ \\

    \theutterance \stepcounter{utterance}  
    & & & \multicolumn{2}{p{0.3\linewidth}}{
        \cellcolor[rgb]{0.9,0.9,0.9}{
            \makecell[{{p{\linewidth}}}]{
                \texttt{\tiny{[GM$|$GM]}}
                {'word': 'kanada', 'assignment': 'team'}
            }
        }
    }
    & & \\ \\

    \theutterance \stepcounter{utterance}  
    & & & \multicolumn{2}{p{0.3\linewidth}}{
        \cellcolor[rgb]{0.9,0.9,0.9}{
            \makecell[{{p{\linewidth}}}]{
                \texttt{\tiny{[GM$|$GM]}}
                {'word': 'indija', 'assignment': 'innocent'}
            }
        }
    }
    & & \\ \\

    \theutterance \stepcounter{utterance}  
    & & & \multicolumn{2}{p{0.3\linewidth}}{
        \cellcolor[rgb]{0.9,0.9,0.9}{
            \makecell[{{p{\linewidth}}}]{
                \texttt{\tiny{[GM$|$GM]}}
                \texttt{indija} \\
            }
        }
    }
    & & \\ \\

    \theutterance \stepcounter{utterance}  
    & & & \multicolumn{2}{p{0.3\linewidth}}{
        \cellcolor[rgb]{0.9,0.9,0.9}{
            \makecell[{{p{\linewidth}}}]{
                \texttt{\tiny{[GM$|$GM]}}
                {'word': 'pajek', 'assignment': 'opponent'}
            }
        }
    }
    & & \\ \\

    \theutterance \stepcounter{utterance}  
    & & & \multicolumn{2}{p{0.3\linewidth}}{
        \cellcolor[rgb]{0.9,0.9,0.9}{
            \makecell[{{p{\linewidth}}}]{
                \texttt{\tiny{[GM$|$GM]}}
                {'hidden': {'team': ['bermuda', 'smrt', 'netopir', 'london', 'miška', 'stadion', 'spojina', 'princesa'], 'innocent': ['nebotičnik', 'trikotnik', 'strežnik', 'helikopter', 'antarktika', 'zob'], 'opponent': ['diamant', 'zaloga', 'krt', 'sesalnik', 'marmor', 'skrinja', 'feniks'], 'assassin': ['vlak']}, 'revealed': {'team': {'team': ['kanada'], 'innocent': ['indija'], 'opponent': [], 'assassin': []}, 'opponent': {'team': [], 'innocent': [], 'opponent': ['pajek'], 'assassin': []}}}
            }
        }
    }
    & & \\ \\

    \theutterance \stepcounter{utterance}  
    & & \multicolumn{4}{p{0.6\linewidth}}{
        \cellcolor[rgb]{0.9,0.9,0.9}{
            \makecell[{{p{\linewidth}}}]{
                \texttt{\tiny{[P1$\langle$GM]}}
                \texttt{Beseda kanada je bila pravilno uganjena. Beseda indija je bila uganjena, vendar je nevinen beseda. Potezo vašega soigralca se je tam končala.} \\
\texttt{Sedaj podajte še en namig, povezan z nekaterimi vašimi preostalimi besedami ekipe, in seznam povezanih ciljnih besed. Ne pozabite začeti namiga z 'NAMIG: ', nato dodati novo vrstico in začeti seznam ciljnih besed, ločenih z vejico, z 'CILJI: '. Opomba: nekatere besede so bile s seznamov odstranjene v primerjavi s prejšnjimi zahtevami.} \\
\\ 
\texttt{Vaše preostale besede ekipe so: bermuda, smrt, netopir, london, miška, stadion, spojina, princesa.} \\
\texttt{Preostale besede nasprotnika so: diamant, zaloga, krt, pajek, sesalnik, marmor, skrinja, feniks.} \\
\texttt{Preostale besede, ki motijo, so: nebotičnik, trikotnik, strežnik, helikopter, antarktika, zob.} \\
\texttt{Preostale smrtonosne besede so: vlak.} \\
            }
        }
    }
    & \\ \\

    \theutterance \stepcounter{utterance}  
    & \multicolumn{4}{p{0.6\linewidth}}{
        \cellcolor[rgb]{0.9,0.9,0.9}{
            \makecell[{{p{\linewidth}}}]{
                \texttt{\tiny{[P1$\rangle$GM]}}
                \texttt{NAMIG: britanija} \\
\texttt{CILJI: london, bermuda} \\
            }
        }
    }
    & & \\ \\

    \theutterance \stepcounter{utterance}  
    & & & \multicolumn{2}{p{0.3\linewidth}}{
        \cellcolor[rgb]{0.9,0.9,0.9}{
            \makecell[{{p{\linewidth}}}]{
                \texttt{\tiny{[GM$|$GM]}}
                \texttt{britanija} \\
            }
        }
    }
    & & \\ \\

    \theutterance \stepcounter{utterance}  
    & & & \multicolumn{2}{p{0.3\linewidth}}{
        \cellcolor[rgb]{0.9,0.9,0.9}{
            \makecell[{{p{\linewidth}}}]{
                \texttt{\tiny{[GM$|$GM]}}
                ['london', 'bermuda']
            }
        }
    }
    & & \\ \\

    \theutterance \stepcounter{utterance}  
    & & & \multicolumn{2}{p{0.3\linewidth}}{
        \cellcolor[rgb]{0.9,0.9,0.9}{
            \makecell[{{p{\linewidth}}}]{
                \texttt{\tiny{[GM$|$GM]}}
                {'word': 'london', 'assignment': 'team'}
            }
        }
    }
    & & \\ \\

    \theutterance \stepcounter{utterance}  
    & & & \multicolumn{2}{p{0.3\linewidth}}{
        \cellcolor[rgb]{0.9,0.9,0.9}{
            \makecell[{{p{\linewidth}}}]{
                \texttt{\tiny{[GM$|$GM]}}
                {'word': 'bermuda', 'assignment': 'team'}
            }
        }
    }
    & & \\ \\

    \theutterance \stepcounter{utterance}  
    & & \multicolumn{4}{p{0.6\linewidth}}{
        \cellcolor[rgb]{0.9,0.9,0.9}{
            \makecell[{{p{\linewidth}}}]{
                \texttt{\tiny{[GM$\rangle$P2]}}
                \texttt{Beseda kanada je bila pravilno uganjena. Beseda indija je bila uganjena, vendar je nevinen beseda. Vaša poteza se je tam končala.} \\
\texttt{Zdaj navedite še en seznam z vejico ločenih besed, najmanj 1 in največ 2 besed z naslednjega seznama besed, ki se najbolj nanašajo ali so najbolj povezane z besedo 'britanija'. Ne pozabite začeti odgovora z UGIBANJE: . Opozorilo: nekatere besede so bile v primerjavi s prejšnjimi zahtevami odstranjene s seznama.} \\
\\ 
\texttt{strežnik, helikopter, netopir, stadion, vlak, sesalnik, bermuda, diamant, nebotičnik, marmor, trikotnik, spojina, krt, zob, smrt, london, miška, skrinja, antarktika, zaloga, feniks, princesa} \\
            }
        }
    }
    & \\ \\

    \theutterance \stepcounter{utterance}  
    & & & \multicolumn{4}{p{0.6\linewidth}}{
        \cellcolor[rgb]{0.9,0.9,0.9}{
            \makecell[{{p{\linewidth}}}]{
                \texttt{\tiny{[GM$\langle$P2]}}
                \texttt{UGIBANJE: london, bermuda} \\
            }
        }
    }
     \\ \\

    \theutterance \stepcounter{utterance}  
    & & & \multicolumn{2}{p{0.3\linewidth}}{
        \cellcolor[rgb]{0.9,0.9,0.9}{
            \makecell[{{p{\linewidth}}}]{
                \texttt{\tiny{[GM$|$GM]}}
                ['london', 'bermuda']
            }
        }
    }
    & & \\ \\

    \theutterance \stepcounter{utterance}  
    & & & \multicolumn{2}{p{0.3\linewidth}}{
        \cellcolor[rgb]{0.9,0.9,0.9}{
            \makecell[{{p{\linewidth}}}]{
                \texttt{\tiny{[GM$|$GM]}}
                {'word': 'london', 'assignment': 'team'}
            }
        }
    }
    & & \\ \\

    \theutterance \stepcounter{utterance}  
    & & & \multicolumn{2}{p{0.3\linewidth}}{
        \cellcolor[rgb]{0.9,0.9,0.9}{
            \makecell[{{p{\linewidth}}}]{
                \texttt{\tiny{[GM$|$GM]}}
                {'word': 'london', 'assignment': 'team'}
            }
        }
    }
    & & \\ \\

    \theutterance \stepcounter{utterance}  
    & & & \multicolumn{2}{p{0.3\linewidth}}{
        \cellcolor[rgb]{0.9,0.9,0.9}{
            \makecell[{{p{\linewidth}}}]{
                \texttt{\tiny{[GM$|$GM]}}
                {'word': 'bermuda', 'assignment': 'team'}
            }
        }
    }
    & & \\ \\

    \theutterance \stepcounter{utterance}  
    & & & \multicolumn{2}{p{0.3\linewidth}}{
        \cellcolor[rgb]{0.9,0.9,0.9}{
            \makecell[{{p{\linewidth}}}]{
                \texttt{\tiny{[GM$|$GM]}}
                {'word': 'bermuda', 'assignment': 'team'}
            }
        }
    }
    & & \\ \\

    \theutterance \stepcounter{utterance}  
    & & & \multicolumn{2}{p{0.3\linewidth}}{
        \cellcolor[rgb]{0.9,0.9,0.9}{
            \makecell[{{p{\linewidth}}}]{
                \texttt{\tiny{[GM$|$GM]}}
                {'word': 'skrinja', 'assignment': 'opponent'}
            }
        }
    }
    & & \\ \\

    \theutterance \stepcounter{utterance}  
    & & & \multicolumn{2}{p{0.3\linewidth}}{
        \cellcolor[rgb]{0.9,0.9,0.9}{
            \makecell[{{p{\linewidth}}}]{
                \texttt{\tiny{[GM$|$GM]}}
                {'hidden': {'team': ['smrt', 'netopir', 'miška', 'stadion', 'spojina', 'princesa'], 'innocent': ['nebotičnik', 'trikotnik', 'strežnik', 'helikopter', 'antarktika', 'zob'], 'opponent': ['diamant', 'zaloga', 'krt', 'sesalnik', 'marmor', 'feniks'], 'assassin': ['vlak']}, 'revealed': {'team': {'team': ['kanada', 'london', 'bermuda'], 'innocent': ['indija'], 'opponent': [], 'assassin': []}, 'opponent': {'team': [], 'innocent': [], 'opponent': ['pajek', 'skrinja'], 'assassin': []}}}
            }
        }
    }
    & & \\ \\

    \theutterance \stepcounter{utterance}  
    & & \multicolumn{4}{p{0.6\linewidth}}{
        \cellcolor[rgb]{0.9,0.9,0.9}{
            \makecell[{{p{\linewidth}}}]{
                \texttt{\tiny{[P1$\langle$GM]}}
                \texttt{Besede london, bermuda so bile pravilno uganjene. Potezo vašega soigralca se je tam končala.} \\
\texttt{Sedaj podajte še en namig, povezan z nekaterimi vašimi preostalimi besedami ekipe, in seznam povezanih ciljnih besed. Ne pozabite začeti namiga z 'NAMIG: ', nato dodati novo vrstico in začeti seznam ciljnih besed, ločenih z vejico, z 'CILJI: '. Opomba: nekatere besede so bile s seznamov odstranjene v primerjavi s prejšnjimi zahtevami.} \\
\\ 
\texttt{Vaše preostale besede ekipe so: smrt, netopir, miška, stadion, spojina, princesa.} \\
\texttt{Preostale besede nasprotnika so: diamant, zaloga, krt, sesalnik, marmor, skrinja, feniks.} \\
\texttt{Preostale besede, ki motijo, so: nebotičnik, trikotnik, strežnik, helikopter, antarktika, zob.} \\
\texttt{Preostale smrtonosne besede so: vlak.} \\
            }
        }
    }
    & \\ \\

    \theutterance \stepcounter{utterance}  
    & \multicolumn{4}{p{0.6\linewidth}}{
        \cellcolor[rgb]{0.9,0.9,0.9}{
            \makecell[{{p{\linewidth}}}]{
                \texttt{\tiny{[P1$\rangle$GM]}}
                \texttt{NAMIG: disney} \\
\texttt{CILJI: princesa, miška} \\
            }
        }
    }
    & & \\ \\

    \theutterance \stepcounter{utterance}  
    & & & \multicolumn{2}{p{0.3\linewidth}}{
        \cellcolor[rgb]{0.9,0.9,0.9}{
            \makecell[{{p{\linewidth}}}]{
                \texttt{\tiny{[GM$|$GM]}}
                \texttt{disney} \\
            }
        }
    }
    & & \\ \\

    \theutterance \stepcounter{utterance}  
    & & & \multicolumn{2}{p{0.3\linewidth}}{
        \cellcolor[rgb]{0.9,0.9,0.9}{
            \makecell[{{p{\linewidth}}}]{
                \texttt{\tiny{[GM$|$GM]}}
                ['princesa', 'miška']
            }
        }
    }
    & & \\ \\

    \theutterance \stepcounter{utterance}  
    & & & \multicolumn{2}{p{0.3\linewidth}}{
        \cellcolor[rgb]{0.9,0.9,0.9}{
            \makecell[{{p{\linewidth}}}]{
                \texttt{\tiny{[GM$|$GM]}}
                {'word': 'princesa', 'assignment': 'team'}
            }
        }
    }
    & & \\ \\

    \theutterance \stepcounter{utterance}  
    & & & \multicolumn{2}{p{0.3\linewidth}}{
        \cellcolor[rgb]{0.9,0.9,0.9}{
            \makecell[{{p{\linewidth}}}]{
                \texttt{\tiny{[GM$|$GM]}}
                {'word': 'miška', 'assignment': 'team'}
            }
        }
    }
    & & \\ \\

    \theutterance \stepcounter{utterance}  
    & & \multicolumn{4}{p{0.6\linewidth}}{
        \cellcolor[rgb]{0.9,0.9,0.9}{
            \makecell[{{p{\linewidth}}}]{
                \texttt{\tiny{[GM$\rangle$P2]}}
                \texttt{Besede london, bermuda so bile pravilno uganjene. Vaša poteza se je tam končala.} \\
\texttt{Zdaj navedite še en seznam z vejico ločenih besed, najmanj 1 in največ 2 besed z naslednjega seznama besed, ki se najbolj nanašajo ali so najbolj povezane z besedo 'disney'. Ne pozabite začeti odgovora z UGIBANJE: . Opozorilo: nekatere besede so bile v primerjavi s prejšnjimi zahtevami odstranjene s seznama.} \\
\\ 
\texttt{strežnik, helikopter, netopir, stadion, vlak, sesalnik, diamant, nebotičnik, marmor, trikotnik, spojina, krt, zob, smrt, miška, antarktika, zaloga, feniks, princesa} \\
            }
        }
    }
    & \\ \\

    \theutterance \stepcounter{utterance}  
    & & & \multicolumn{4}{p{0.6\linewidth}}{
        \cellcolor[rgb]{0.9,0.9,0.9}{
            \makecell[{{p{\linewidth}}}]{
                \texttt{\tiny{[GM$\langle$P2]}}
                \texttt{UGIBANJE: princesa, miška} \\
            }
        }
    }
     \\ \\

    \theutterance \stepcounter{utterance}  
    & & & \multicolumn{2}{p{0.3\linewidth}}{
        \cellcolor[rgb]{0.9,0.9,0.9}{
            \makecell[{{p{\linewidth}}}]{
                \texttt{\tiny{[GM$|$GM]}}
                ['princesa', 'miška']
            }
        }
    }
    & & \\ \\

    \theutterance \stepcounter{utterance}  
    & & & \multicolumn{2}{p{0.3\linewidth}}{
        \cellcolor[rgb]{0.9,0.9,0.9}{
            \makecell[{{p{\linewidth}}}]{
                \texttt{\tiny{[GM$|$GM]}}
                {'word': 'princesa', 'assignment': 'team'}
            }
        }
    }
    & & \\ \\

    \theutterance \stepcounter{utterance}  
    & & & \multicolumn{2}{p{0.3\linewidth}}{
        \cellcolor[rgb]{0.9,0.9,0.9}{
            \makecell[{{p{\linewidth}}}]{
                \texttt{\tiny{[GM$|$GM]}}
                {'word': 'princesa', 'assignment': 'team'}
            }
        }
    }
    & & \\ \\

    \theutterance \stepcounter{utterance}  
    & & & \multicolumn{2}{p{0.3\linewidth}}{
        \cellcolor[rgb]{0.9,0.9,0.9}{
            \makecell[{{p{\linewidth}}}]{
                \texttt{\tiny{[GM$|$GM]}}
                {'word': 'miška', 'assignment': 'team'}
            }
        }
    }
    & & \\ \\

    \theutterance \stepcounter{utterance}  
    & & & \multicolumn{2}{p{0.3\linewidth}}{
        \cellcolor[rgb]{0.9,0.9,0.9}{
            \makecell[{{p{\linewidth}}}]{
                \texttt{\tiny{[GM$|$GM]}}
                {'word': 'miška', 'assignment': 'team'}
            }
        }
    }
    & & \\ \\

    \theutterance \stepcounter{utterance}  
    & & & \multicolumn{2}{p{0.3\linewidth}}{
        \cellcolor[rgb]{0.9,0.9,0.9}{
            \makecell[{{p{\linewidth}}}]{
                \texttt{\tiny{[GM$|$GM]}}
                {'word': 'krt', 'assignment': 'opponent'}
            }
        }
    }
    & & \\ \\

    \theutterance \stepcounter{utterance}  
    & & & \multicolumn{2}{p{0.3\linewidth}}{
        \cellcolor[rgb]{0.9,0.9,0.9}{
            \makecell[{{p{\linewidth}}}]{
                \texttt{\tiny{[GM$|$GM]}}
                {'hidden': {'team': ['smrt', 'netopir', 'stadion', 'spojina'], 'innocent': ['nebotičnik', 'trikotnik', 'strežnik', 'helikopter', 'antarktika', 'zob'], 'opponent': ['diamant', 'zaloga', 'sesalnik', 'marmor', 'feniks'], 'assassin': ['vlak']}, 'revealed': {'team': {'team': ['kanada', 'london', 'bermuda', 'princesa', 'miška'], 'innocent': ['indija'], 'opponent': [], 'assassin': []}, 'opponent': {'team': [], 'innocent': [], 'opponent': ['pajek', 'skrinja', 'krt'], 'assassin': []}}}
            }
        }
    }
    & & \\ \\

    \theutterance \stepcounter{utterance}  
    & & \multicolumn{4}{p{0.6\linewidth}}{
        \cellcolor[rgb]{0.9,0.9,0.9}{
            \makecell[{{p{\linewidth}}}]{
                \texttt{\tiny{[P1$\langle$GM]}}
                \texttt{Besede princesa, miška so bile pravilno uganjene. Potezo vašega soigralca se je tam končala.} \\
\texttt{Sedaj podajte še en namig, povezan z nekaterimi vašimi preostalimi besedami ekipe, in seznam povezanih ciljnih besed. Ne pozabite začeti namiga z 'NAMIG: ', nato dodati novo vrstico in začeti seznam ciljnih besed, ločenih z vejico, z 'CILJI: '. Opomba: nekatere besede so bile s seznamov odstranjene v primerjavi s prejšnjimi zahtevami.} \\
\\ 
\texttt{Vaše preostale besede ekipe so: smrt, netopir, stadion, spojina.} \\
\texttt{Preostale besede nasprotnika so: diamant, zaloga, krt, sesalnik, marmor, feniks.} \\
\texttt{Preostale besede, ki motijo, so: nebotičnik, trikotnik, strežnik, helikopter, antarktika, zob.} \\
\texttt{Preostale smrtonosne besede so: vlak.} \\
            }
        }
    }
    & \\ \\

    \theutterance \stepcounter{utterance}  
    & \multicolumn{4}{p{0.6\linewidth}}{
        \cellcolor[rgb]{0.9,0.9,0.9}{
            \makecell[{{p{\linewidth}}}]{
                \texttt{\tiny{[P1$\rangle$GM]}}
                \texttt{NAMIG: vampir} \\
\texttt{CILJI: netopir, smrt} \\
            }
        }
    }
    & & \\ \\

    \theutterance \stepcounter{utterance}  
    & & & \multicolumn{2}{p{0.3\linewidth}}{
        \cellcolor[rgb]{0.9,0.9,0.9}{
            \makecell[{{p{\linewidth}}}]{
                \texttt{\tiny{[GM$|$GM]}}
                \texttt{vampir} \\
            }
        }
    }
    & & \\ \\

    \theutterance \stepcounter{utterance}  
    & & & \multicolumn{2}{p{0.3\linewidth}}{
        \cellcolor[rgb]{0.9,0.9,0.9}{
            \makecell[{{p{\linewidth}}}]{
                \texttt{\tiny{[GM$|$GM]}}
                ['netopir', 'smrt']
            }
        }
    }
    & & \\ \\

    \theutterance \stepcounter{utterance}  
    & & & \multicolumn{2}{p{0.3\linewidth}}{
        \cellcolor[rgb]{0.9,0.9,0.9}{
            \makecell[{{p{\linewidth}}}]{
                \texttt{\tiny{[GM$|$GM]}}
                {'word': 'netopir', 'assignment': 'team'}
            }
        }
    }
    & & \\ \\

    \theutterance \stepcounter{utterance}  
    & & & \multicolumn{2}{p{0.3\linewidth}}{
        \cellcolor[rgb]{0.9,0.9,0.9}{
            \makecell[{{p{\linewidth}}}]{
                \texttt{\tiny{[GM$|$GM]}}
                {'word': 'smrt', 'assignment': 'team'}
            }
        }
    }
    & & \\ \\

    \theutterance \stepcounter{utterance}  
    & & \multicolumn{4}{p{0.6\linewidth}}{
        \cellcolor[rgb]{0.9,0.9,0.9}{
            \makecell[{{p{\linewidth}}}]{
                \texttt{\tiny{[GM$\rangle$P2]}}
                \texttt{Besede princesa, miška so bile pravilno uganjene. Vaša poteza se je tam končala.} \\
\texttt{Zdaj navedite še en seznam z vejico ločenih besed, najmanj 1 in največ 2 besed z naslednjega seznama besed, ki se najbolj nanašajo ali so najbolj povezane z besedo 'vampir'. Ne pozabite začeti odgovora z UGIBANJE: . Opozorilo: nekatere besede so bile v primerjavi s prejšnjimi zahtevami odstranjene s seznama.} \\
\\ 
\texttt{strežnik, helikopter, netopir, stadion, vlak, sesalnik, diamant, nebotičnik, marmor, trikotnik, spojina, zob, smrt, antarktika, zaloga, feniks} \\
            }
        }
    }
    & \\ \\

    \theutterance \stepcounter{utterance}  
    & & & \multicolumn{4}{p{0.6\linewidth}}{
        \cellcolor[rgb]{0.9,0.9,0.9}{
            \makecell[{{p{\linewidth}}}]{
                \texttt{\tiny{[GM$\langle$P2]}}
                \texttt{UGIBANJE: netopir, zob} \\
            }
        }
    }
     \\ \\

    \theutterance \stepcounter{utterance}  
    & & & \multicolumn{2}{p{0.3\linewidth}}{
        \cellcolor[rgb]{0.9,0.9,0.9}{
            \makecell[{{p{\linewidth}}}]{
                \texttt{\tiny{[GM$|$GM]}}
                ['netopir', 'zob']
            }
        }
    }
    & & \\ \\

    \theutterance \stepcounter{utterance}  
    & & & \multicolumn{2}{p{0.3\linewidth}}{
        \cellcolor[rgb]{0.9,0.9,0.9}{
            \makecell[{{p{\linewidth}}}]{
                \texttt{\tiny{[GM$|$GM]}}
                {'word': 'netopir', 'assignment': 'team'}
            }
        }
    }
    & & \\ \\

    \theutterance \stepcounter{utterance}  
    & & & \multicolumn{2}{p{0.3\linewidth}}{
        \cellcolor[rgb]{0.9,0.9,0.9}{
            \makecell[{{p{\linewidth}}}]{
                \texttt{\tiny{[GM$|$GM]}}
                {'word': 'netopir', 'assignment': 'team'}
            }
        }
    }
    & & \\ \\

    \theutterance \stepcounter{utterance}  
    & & & \multicolumn{2}{p{0.3\linewidth}}{
        \cellcolor[rgb]{0.9,0.9,0.9}{
            \makecell[{{p{\linewidth}}}]{
                \texttt{\tiny{[GM$|$GM]}}
                {'word': 'zob', 'assignment': 'innocent'}
            }
        }
    }
    & & \\ \\

    \theutterance \stepcounter{utterance}  
    & & & \multicolumn{2}{p{0.3\linewidth}}{
        \cellcolor[rgb]{0.9,0.9,0.9}{
            \makecell[{{p{\linewidth}}}]{
                \texttt{\tiny{[GM$|$GM]}}
                \texttt{zob} \\
            }
        }
    }
    & & \\ \\

    \theutterance \stepcounter{utterance}  
    & & & \multicolumn{2}{p{0.3\linewidth}}{
        \cellcolor[rgb]{0.9,0.9,0.9}{
            \makecell[{{p{\linewidth}}}]{
                \texttt{\tiny{[GM$|$GM]}}
                {'word': 'diamant', 'assignment': 'opponent'}
            }
        }
    }
    & & \\ \\

    \theutterance \stepcounter{utterance}  
    & & & \multicolumn{2}{p{0.3\linewidth}}{
        \cellcolor[rgb]{0.9,0.9,0.9}{
            \makecell[{{p{\linewidth}}}]{
                \texttt{\tiny{[GM$|$GM]}}
                {'hidden': {'team': ['smrt', 'stadion', 'spojina'], 'innocent': ['nebotičnik', 'trikotnik', 'strežnik', 'helikopter', 'antarktika'], 'opponent': ['zaloga', 'sesalnik', 'marmor', 'feniks'], 'assassin': ['vlak']}, 'revealed': {'team': {'team': ['kanada', 'london', 'bermuda', 'princesa', 'miška', 'netopir'], 'innocent': ['indija', 'zob'], 'opponent': [], 'assassin': []}, 'opponent': {'team': [], 'innocent': [], 'opponent': ['pajek', 'skrinja', 'krt', 'diamant'], 'assassin': []}}}
            }
        }
    }
    & & \\ \\

    \theutterance \stepcounter{utterance}  
    & & \multicolumn{4}{p{0.6\linewidth}}{
        \cellcolor[rgb]{0.9,0.9,0.9}{
            \makecell[{{p{\linewidth}}}]{
                \texttt{\tiny{[P1$\langle$GM]}}
                \texttt{Beseda netopir je bila pravilno uganjena. Beseda zob je bila uganjena, vendar je nevinen beseda. Potezo vašega soigralca se je tam končala.} \\
\texttt{Sedaj podajte še en namig, povezan z nekaterimi vašimi preostalimi besedami ekipe, in seznam povezanih ciljnih besed. Ne pozabite začeti namiga z 'NAMIG: ', nato dodati novo vrstico in začeti seznam ciljnih besed, ločenih z vejico, z 'CILJI: '. Opomba: nekatere besede so bile s seznamov odstranjene v primerjavi s prejšnjimi zahtevami.} \\
\\ 
\texttt{Vaše preostale besede ekipe so: smrt, stadion, spojina.} \\
\texttt{Preostale besede nasprotnika so: diamant, zaloga, sesalnik, marmor, feniks.} \\
\texttt{Preostale besede, ki motijo, so: nebotičnik, trikotnik, strežnik, helikopter, antarktika.} \\
\texttt{Preostale smrtonosne besede so: vlak.} \\
            }
        }
    }
    & \\ \\

    \theutterance \stepcounter{utterance}  
    & \multicolumn{4}{p{0.6\linewidth}}{
        \cellcolor[rgb]{0.9,0.9,0.9}{
            \makecell[{{p{\linewidth}}}]{
                \texttt{\tiny{[P1$\rangle$GM]}}
                \texttt{NAMIG: kemija} \\
\texttt{CILJI: spojina} \\
            }
        }
    }
    & & \\ \\

    \theutterance \stepcounter{utterance}  
    & & & \multicolumn{2}{p{0.3\linewidth}}{
        \cellcolor[rgb]{0.9,0.9,0.9}{
            \makecell[{{p{\linewidth}}}]{
                \texttt{\tiny{[GM$|$GM]}}
                \texttt{kemija} \\
            }
        }
    }
    & & \\ \\

    \theutterance \stepcounter{utterance}  
    & & & \multicolumn{2}{p{0.3\linewidth}}{
        \cellcolor[rgb]{0.9,0.9,0.9}{
            \makecell[{{p{\linewidth}}}]{
                \texttt{\tiny{[GM$|$GM]}}
                ['spojina']
            }
        }
    }
    & & \\ \\

    \theutterance \stepcounter{utterance}  
    & & & \multicolumn{2}{p{0.3\linewidth}}{
        \cellcolor[rgb]{0.9,0.9,0.9}{
            \makecell[{{p{\linewidth}}}]{
                \texttt{\tiny{[GM$|$GM]}}
                {'word': 'spojina', 'assignment': 'team'}
            }
        }
    }
    & & \\ \\

    \theutterance \stepcounter{utterance}  
    & & \multicolumn{4}{p{0.6\linewidth}}{
        \cellcolor[rgb]{0.9,0.9,0.9}{
            \makecell[{{p{\linewidth}}}]{
                \texttt{\tiny{[GM$\rangle$P2]}}
                \texttt{Beseda netopir je bila pravilno uganjena. Beseda zob je bila uganjena, vendar je nevinen beseda. Vaša poteza se je tam končala.} \\
\texttt{Zdaj navedite še en seznam z vejico ločenih besed, najmanj 1 in največ 1 besed z naslednjega seznama besed, ki se najbolj nanašajo ali so najbolj povezane z besedo 'kemija'. Ne pozabite začeti odgovora z UGIBANJE: . Opozorilo: nekatere besede so bile v primerjavi s prejšnjimi zahtevami odstranjene s seznama.} \\
\\ 
\texttt{strežnik, helikopter, stadion, vlak, sesalnik, nebotičnik, marmor, trikotnik, spojina, smrt, antarktika, zaloga, feniks} \\
            }
        }
    }
    & \\ \\

    \theutterance \stepcounter{utterance}  
    & & & \multicolumn{4}{p{0.6\linewidth}}{
        \cellcolor[rgb]{0.9,0.9,0.9}{
            \makecell[{{p{\linewidth}}}]{
                \texttt{\tiny{[GM$\langle$P2]}}
                \texttt{UGIBANJE: spojina} \\
            }
        }
    }
     \\ \\

    \theutterance \stepcounter{utterance}  
    & & & \multicolumn{2}{p{0.3\linewidth}}{
        \cellcolor[rgb]{0.9,0.9,0.9}{
            \makecell[{{p{\linewidth}}}]{
                \texttt{\tiny{[GM$|$GM]}}
                ['spojina']
            }
        }
    }
    & & \\ \\

    \theutterance \stepcounter{utterance}  
    & & & \multicolumn{2}{p{0.3\linewidth}}{
        \cellcolor[rgb]{0.9,0.9,0.9}{
            \makecell[{{p{\linewidth}}}]{
                \texttt{\tiny{[GM$|$GM]}}
                {'word': 'spojina', 'assignment': 'team'}
            }
        }
    }
    & & \\ \\

    \theutterance \stepcounter{utterance}  
    & & & \multicolumn{2}{p{0.3\linewidth}}{
        \cellcolor[rgb]{0.9,0.9,0.9}{
            \makecell[{{p{\linewidth}}}]{
                \texttt{\tiny{[GM$|$GM]}}
                {'word': 'spojina', 'assignment': 'team'}
            }
        }
    }
    & & \\ \\

    \theutterance \stepcounter{utterance}  
    & & & \multicolumn{2}{p{0.3\linewidth}}{
        \cellcolor[rgb]{0.9,0.9,0.9}{
            \makecell[{{p{\linewidth}}}]{
                \texttt{\tiny{[GM$|$GM]}}
                {'word': 'sesalnik', 'assignment': 'opponent'}
            }
        }
    }
    & & \\ \\

    \theutterance \stepcounter{utterance}  
    & & & \multicolumn{2}{p{0.3\linewidth}}{
        \cellcolor[rgb]{0.9,0.9,0.9}{
            \makecell[{{p{\linewidth}}}]{
                \texttt{\tiny{[GM$|$GM]}}
                {'hidden': {'team': ['smrt', 'stadion'], 'innocent': ['nebotičnik', 'trikotnik', 'strežnik', 'helikopter', 'antarktika'], 'opponent': ['zaloga', 'marmor', 'feniks'], 'assassin': ['vlak']}, 'revealed': {'team': {'team': ['kanada', 'london', 'bermuda', 'princesa', 'miška', 'netopir', 'spojina'], 'innocent': ['indija', 'zob'], 'opponent': [], 'assassin': []}, 'opponent': {'team': [], 'innocent': [], 'opponent': ['pajek', 'skrinja', 'krt', 'diamant', 'sesalnik'], 'assassin': []}}}
            }
        }
    }
    & & \\ \\

    \theutterance \stepcounter{utterance}  
    & & \multicolumn{4}{p{0.6\linewidth}}{
        \cellcolor[rgb]{0.9,0.9,0.9}{
            \makecell[{{p{\linewidth}}}]{
                \texttt{\tiny{[P1$\langle$GM]}}
                \texttt{Beseda spojina je bila pravilno uganjena. Potezo vašega soigralca se je tam končala.} \\
\texttt{Sedaj podajte še en namig, povezan z nekaterimi vašimi preostalimi besedami ekipe, in seznam povezanih ciljnih besed. Ne pozabite začeti namiga z 'NAMIG: ', nato dodati novo vrstico in začeti seznam ciljnih besed, ločenih z vejico, z 'CILJI: '. Opomba: nekatere besede so bile s seznamov odstranjene v primerjavi s prejšnjimi zahtevami.} \\
\\ 
\texttt{Vaše preostale besede ekipe so: smrt, stadion.} \\
\texttt{Preostale besede nasprotnika so: zaloga, sesalnik, marmor, feniks.} \\
\texttt{Preostale besede, ki motijo, so: nebotičnik, trikotnik, strežnik, helikopter, antarktika.} \\
\texttt{Preostale smrtonosne besede so: vlak.} \\
            }
        }
    }
    & \\ \\

    \theutterance \stepcounter{utterance}  
    & \multicolumn{4}{p{0.6\linewidth}}{
        \cellcolor[rgb]{0.9,0.9,0.9}{
            \makecell[{{p{\linewidth}}}]{
                \texttt{\tiny{[P1$\rangle$GM]}}
                \texttt{NAMIG: gladiator} \\
\texttt{CILJI: stadion, smrt} \\
            }
        }
    }
    & & \\ \\

    \theutterance \stepcounter{utterance}  
    & & & \multicolumn{2}{p{0.3\linewidth}}{
        \cellcolor[rgb]{0.9,0.9,0.9}{
            \makecell[{{p{\linewidth}}}]{
                \texttt{\tiny{[GM$|$GM]}}
                \texttt{gladiator} \\
            }
        }
    }
    & & \\ \\

    \theutterance \stepcounter{utterance}  
    & & & \multicolumn{2}{p{0.3\linewidth}}{
        \cellcolor[rgb]{0.9,0.9,0.9}{
            \makecell[{{p{\linewidth}}}]{
                \texttt{\tiny{[GM$|$GM]}}
                ['stadion', 'smrt']
            }
        }
    }
    & & \\ \\

    \theutterance \stepcounter{utterance}  
    & & & \multicolumn{2}{p{0.3\linewidth}}{
        \cellcolor[rgb]{0.9,0.9,0.9}{
            \makecell[{{p{\linewidth}}}]{
                \texttt{\tiny{[GM$|$GM]}}
                {'word': 'stadion', 'assignment': 'team'}
            }
        }
    }
    & & \\ \\

    \theutterance \stepcounter{utterance}  
    & & & \multicolumn{2}{p{0.3\linewidth}}{
        \cellcolor[rgb]{0.9,0.9,0.9}{
            \makecell[{{p{\linewidth}}}]{
                \texttt{\tiny{[GM$|$GM]}}
                {'word': 'smrt', 'assignment': 'team'}
            }
        }
    }
    & & \\ \\

    \theutterance \stepcounter{utterance}  
    & & \multicolumn{4}{p{0.6\linewidth}}{
        \cellcolor[rgb]{0.9,0.9,0.9}{
            \makecell[{{p{\linewidth}}}]{
                \texttt{\tiny{[GM$\rangle$P2]}}
                \texttt{Beseda spojina je bila pravilno uganjena. Vaša poteza se je tam končala.} \\
\texttt{Zdaj navedite še en seznam z vejico ločenih besed, najmanj 1 in največ 2 besed z naslednjega seznama besed, ki se najbolj nanašajo ali so najbolj povezane z besedo 'gladiator'. Ne pozabite začeti odgovora z UGIBANJE: . Opozorilo: nekatere besede so bile v primerjavi s prejšnjimi zahtevami odstranjene s seznama.} \\
\\ 
\texttt{strežnik, helikopter, stadion, vlak, nebotičnik, marmor, trikotnik, smrt, antarktika, zaloga, feniks} \\
            }
        }
    }
    & \\ \\

    \theutterance \stepcounter{utterance}  
    & & & \multicolumn{4}{p{0.6\linewidth}}{
        \cellcolor[rgb]{0.9,0.9,0.9}{
            \makecell[{{p{\linewidth}}}]{
                \texttt{\tiny{[GM$\langle$P2]}}
                \texttt{UGIBANJE: stadion, smrt} \\
            }
        }
    }
     \\ \\

    \theutterance \stepcounter{utterance}  
    & & & \multicolumn{2}{p{0.3\linewidth}}{
        \cellcolor[rgb]{0.9,0.9,0.9}{
            \makecell[{{p{\linewidth}}}]{
                \texttt{\tiny{[GM$|$GM]}}
                ['stadion', 'smrt']
            }
        }
    }
    & & \\ \\

    \theutterance \stepcounter{utterance}  
    & & & \multicolumn{2}{p{0.3\linewidth}}{
        \cellcolor[rgb]{0.9,0.9,0.9}{
            \makecell[{{p{\linewidth}}}]{
                \texttt{\tiny{[GM$|$GM]}}
                {'word': 'stadion', 'assignment': 'team'}
            }
        }
    }
    & & \\ \\

    \theutterance \stepcounter{utterance}  
    & & & \multicolumn{2}{p{0.3\linewidth}}{
        \cellcolor[rgb]{0.9,0.9,0.9}{
            \makecell[{{p{\linewidth}}}]{
                \texttt{\tiny{[GM$|$GM]}}
                {'word': 'stadion', 'assignment': 'team'}
            }
        }
    }
    & & \\ \\

    \theutterance \stepcounter{utterance}  
    & & & \multicolumn{2}{p{0.3\linewidth}}{
        \cellcolor[rgb]{0.9,0.9,0.9}{
            \makecell[{{p{\linewidth}}}]{
                \texttt{\tiny{[GM$|$GM]}}
                {'word': 'smrt', 'assignment': 'team'}
            }
        }
    }
    & & \\ \\

    \theutterance \stepcounter{utterance}  
    & & & \multicolumn{2}{p{0.3\linewidth}}{
        \cellcolor[rgb]{0.9,0.9,0.9}{
            \makecell[{{p{\linewidth}}}]{
                \texttt{\tiny{[GM$|$GM]}}
                {'word': 'smrt', 'assignment': 'team'}
            }
        }
    }
    & & \\ \\

    \theutterance \stepcounter{utterance}  
    & & & \multicolumn{2}{p{0.3\linewidth}}{
        \cellcolor[rgb]{0.9,0.9,0.9}{
            \makecell[{{p{\linewidth}}}]{
                \texttt{\tiny{[GM$|$GM]}}
                \texttt{team has won} \\
            }
        }
    }
    & & \\ \\

\end{supertabular}
}

\end{document}
